% Original pages: 97-104
% Figures: none.
% Uncertainty log: none.

\subsection{Grothendieck groups}

\noindent\textbf{26.} Let \(A\) be a Noetherian ring and let \(F(A)\) denote the set of all isomorphism classes of finitely generated \(A\)-modules. Let \(C\) be the free abelian group generated by \(F(A)\). With each short exact sequence \(0\to M'\to M\to M''\to0\) of finitely generated \(A\)-modules we associate the element \((M')-(M)+(M'')\) of \(C\), where \((M)\) is the isomorphism class of \(M\), etc. Let \(D\) be the subgroup of \(C\) generated by these elements, for all short exact sequences. The quotient group \(C/D\) is called the \emph{Grothendieck group} of \(A\), and is denoted by \(K(A)\). If \(M\) is a finitely generated \(A\)-module, let \(\gamma(M)\), or \(\gamma_A(M)\), denote the image of \((M)\) in \(K(A)\).

i) Show that \(K(A)\) has the following universal property: for each additive function \(\lambda\) on the class of finitely generated \(A\)-modules, with values in an abelian group \(G\), there exists a unique homomorphism \(\lambda_0:K(A)\to G\) such that \(\lambda(M)=\lambda_0(\gamma(M))\) for all \(M\).

ii) Show that \(K(A)\) is generated by the elements \(\gamma(A/\mathfrak p)\), where \(\mathfrak p\) is a prime ideal of \(A\). [Use Exercise 18.]

iii) If \(A\) is a field, or more generally if \(A\) is a principal ideal domain, then \(K(A)\cong Z\).

iv) Let \(f:A\to B\) be a \emph{finite} ring homomorphism. Show that restriction of scalars gives rise to a homomorphism \(f_!:K(B)\to K(A)\) such that \(f_!(\gamma_B(N))=\gamma_A(N)\) for a \(B\)-module \(N\). If \(g:B\to C\) is another finite ring homomorphism, show that \((g\circ f)_!=f_!\circ g_!\).

\noindent\textbf{27.} Let \(A\) be a Noetherian ring and let \(F_1(A)\) be the set of all isomorphism classes of finitely generated \emph{flat} \(A\)-modules. Repeating the construction of Exercise 26 we obtain a group \(K_1(A)\). Let \(\gamma_1(M)\) denote the image of \((M)\) in \(K_1(A)\).

i) Show that tensor product of modules over \(A\) induces a commutative ring structure on \(K_1(A)\), such that \(\gamma_1(M)\cdot\gamma_1(N)=\gamma_1(M\otimes N)\). The identity element of this ring is \(\gamma_1(A)\).

ii) Show that tensor product induces a \(K_1(A)\)-module structure on the group \(K(A)\), such that \(\gamma_1(M)\cdot\gamma(N)=\gamma(M\otimes N)\).

iii) If \(A\) is a (Noetherian) local ring, then \(K_1(A)\cong Z\).

iv) Let \(f:A\to B\) be a ring homomorphism, \(B\) being Noetherian. Show that extension of scalars gives rise to a ring homomorphism \(f^!:K_1(A)\to K_1(B)\) such that \(f^!(\gamma_1(M))=\gamma_1(B\otimes_A M)\). [If \(M\) is flat and finitely generated over \(A\), then \(B\otimes_A M\) is flat and finitely generated over \(B\).] If \(g:B\to C\) is another ring homomorphism (with \(C\) Noetherian), then \((f\circ g)^!=f^!\circ g^!\).

v) If \(f:A\to B\) is a finite ring homomorphism then
\[
f_!(f^!(x)y)=xf_!(y)
\]
for \(x\in K_1(A)\), \(y\in K(B)\). In other words, regarding \(K(B)\) as a \(K_1(A)\)-module by restriction of scalars, the homomorphism \(f_!\) is a \(K_1(A)\)-module homomorphism.

\noindent\emph{Remark.} Since \(F_1(A)\) is a subset of \(F(A)\) we have a group homomorphism \(\varepsilon:K_1(A)\to K(A)\), given by \(\varepsilon(\gamma_1(M))=\gamma(M)\). If the ring \(A\) is finite-dimensional and \emph{regular}, i.e., if all its local rings \(A_\mathfrak p\) are regular (Chapter 11) it can be shown that \(\varepsilon\) is an isomorphism.

\section{Artin Rings}

An \emph{Artin ring} is one which satisfies the d.c.c. (or equivalently the minimal condition) on ideals.

The apparent symmetry with Noetherian rings is however misleading. In fact we will show that an Artin ring is necessarily Noetherian and of a very special kind. In a sense an Artin ring is the simplest kind of ring after a field, and we study them not because of their generality but because of their simplicity.

\noindent\textbf{Proposition 8.1.\amtarget{8.1}} \emph{In an Artin ring \(A\) every prime ideal is maximal.}

\noindent\emph{Proof.} Let \(\mathfrak p\) be a prime ideal of \(A\). Then \(B=A/\mathfrak p\) is an Artinian integral domain. Let \(x\in B\), \(x\neq0\). By the d.c.c. we have \((x^n)=(x^{n+1})\) for some \(n\), hence \(x^n=x^{n+1}y\) for some \(y\in B\). Since \(B\) is an integral domain and \(x\neq0\), it follows that we may cancel \(x^n\), hence \(xy=1\). Hence \(x\) has an inverse in \(B\), and therefore \(B\) is a field, so that \(\mathfrak p\) is a maximal ideal. \proofbox

\noindent\textbf{Corollary 8.2.\amtarget{8.2}} \emph{In an Artin ring the nilradical is equal to the Jacobson radical.} \proofbox

\noindent\textbf{Proposition 8.3.\amtarget{8.3}} \emph{An Artin ring has only a finite number of maximal ideals.}

\noindent\emph{Proof.} Consider the set of all finite intersections \(\mathfrak m_1\cap\cdots\cap\mathfrak m_r\), where the \(\mathfrak m_i\) are maximal ideals. This set has a minimal element, say \(\mathfrak m_1\cap\cdots\cap\mathfrak m_n\); hence for any maximal ideal \(\mathfrak m\) we have \(\mathfrak m\cap\mathfrak m_1\cap\cdots\cap\mathfrak m_n=\mathfrak m_1\cap\cdots\cap\mathfrak m_n\), and therefore \(\mathfrak m\supseteq\mathfrak m_1\cap\cdots\cap\mathfrak m_n\). By \amref{1.11} \(\mathfrak m\supseteq\mathfrak m_i\) for some \(i\), hence \(\mathfrak m=\mathfrak m_i\) since \(\mathfrak m_i\) is maximal. \proofbox

\noindent\textbf{Proposition 8.4.\amtarget{8.4}} \emph{In an Artin ring the nilradical \(\mathfrak N\) is nilpotent.}

\noindent\emph{Proof.} By d.c.c. we have \(\mathfrak N^k=\mathfrak N^{k+1}=\cdots=\mathfrak a\) say, for some \(k>0\). Suppose \(\mathfrak a\neq0\), and let \(\Sigma\) denote the set of all ideals \(\mathfrak b\) such that \(\mathfrak a\mathfrak b\neq0\). Then \(\Sigma\) is not empty, since \(\mathfrak a\in\Sigma\). Let \(\mathfrak c\) be a minimal element of \(\Sigma\); then there exists \(x\in\mathfrak c\) such that \(x\mathfrak a\neq0\); we have \((x)\subseteq\mathfrak c\), hence \((x)=\mathfrak c\) by the minimality of \(\mathfrak c\). But \((x\mathfrak a)\mathfrak a=x\mathfrak a^2=x\mathfrak a\neq0\), and \(x\mathfrak a\subseteq(x)\), hence \(x\mathfrak a=(x)\) (again by minimality). Hence \(x=xy\) for some \(y\in\mathfrak a\), and therefore \(x=xy=xy^2=\cdots=xy^n=\cdots\). But \(y\in\mathfrak a=\mathfrak N^k\subseteq\mathfrak N\), hence \(y\) is nilpotent and therefore \(x=xy^n=0\). This contradicts the choice of \(x\), therefore \(\mathfrak a=0\). \proofbox

By a \emph{chain of prime ideals} of a ring \(A\) we mean a finite strictly increasing sequence \(\mathfrak p_0\subset\mathfrak p_1\subset\cdots\subset\mathfrak p_n\); the \emph{length} of the chain is \(n\). We define the \emph{dimension} of \(A\) to be the supremum of the lengths of all chains of prime ideals in \(A\): it is an integer \(\geq0\), or \(+\infty\) (assuming \(A\neq0\)). A field has dimension 0; the ring \(Z\) has dimension 1.

\noindent\textbf{Theorem 8.5.\amtarget{8.5}} \emph{A ring \(A\) is Artin \(\Longleftrightarrow A\) is Noetherian and \(\dim A=0\).}

\noindent\emph{Proof.} \(\Rightarrow\): By \amref{8.1} we have \(\dim A=0\). Let \(\mathfrak m_i\) \((1\leq i\leq n)\) be the distinct maximal ideals of \(A\) \amref{8.3}. Then \(\prod_{i=1}^n\mathfrak m_i^k\subseteq(\bigcap_{i=1}^n\mathfrak m_i)^k=\mathfrak N^k=0\). Hence by \amref{6.11} \(A\) is Noetherian.

\(\Leftarrow\): Since the zero ideal has a primary decomposition \amref{7.13}, \(A\) has only a finite number of minimal prime ideals, and these are all maximal since \(\dim A=0\). Hence \(\mathfrak N=\bigcap_{i=1}^n\mathfrak m_i\) say; we have \(\mathfrak N^k=0\) by \amref{7.15}, hence \(\prod_{i=1}^n\mathfrak m_i^k=0\) as in the previous part of the proof. Hence by \amref{6.11} \(A\) is an Artin ring. \proofbox

If \(A\) is an Artin local ring with maximal ideal \(\mathfrak m\), then \(\mathfrak m\) is the only prime ideal of \(A\) and therefore \(\mathfrak m\) is the nilradical of \(A\). Hence every element of \(\mathfrak m\) is nilpotent, and \(\mathfrak m\) itself is nilpotent. Every element of \(A\) is either a unit or is nilpotent. An example of such a ring is \(Z/(p^n)\), where \(p\) is prime and \(n\geq1\).

\noindent\textbf{Proposition 8.6.\amtarget{8.6}} \emph{Let \(A\) be a Noetherian local ring, \(\mathfrak m\) its maximal ideal. Then exactly one of the following two statements is true:}

\emph{i) \(\mathfrak m^n\neq\mathfrak m^{n+1}\) for all \(n\);}

\emph{ii) \(\mathfrak m^n=0\) for some \(n\), in which case \(A\) is an Artin local ring.}

\noindent\emph{Proof.} Suppose \(\mathfrak m^n=\mathfrak m^{n+1}\) for some \(n\). By Nakayama's lemma \amref{2.6} we have \(\mathfrak m^n=0\). Let \(\mathfrak p\) be any prime ideal of \(A\). Then \(\mathfrak m^n\subseteq\mathfrak p\), hence (taking radicals) \(\mathfrak m=\mathfrak p\). Hence \(\mathfrak m\) is the only prime ideal of \(A\) and therefore \(A\) is Artinian. \proofbox

\noindent\textbf{Theorem 8.7.\amtarget{8.7} (structure theorem for Artin rings).} \emph{An Artin ring \(A\) is uniquely (up to isomorphism) a finite direct product of Artin local rings.}

\noindent\emph{Proof.} Let \(\mathfrak m_i\) \((1\leq i\leq n)\) be the distinct maximal ideals of \(A\). From the proof of \amref{8.5} we have \(\prod_{i=1}^n\mathfrak m_i^k=0\) for some \(k>0\). By \amref{1.16} the ideals \(\mathfrak m_i^k\) are coprime in pairs, hence \(\bigcap\mathfrak m_i^k=\prod\mathfrak m_i^k\) by \amref{1.10}. Consequently by \amref{1.10} again the natural mapping \(A\to\prod_{i=1}^n(A/\mathfrak m_i^k)\) is an isomorphism. Each \(A/\mathfrak m_i^k\) is an Artin local ring, hence \(A\) is a direct product of Artin local rings.

Conversely, suppose \(A\cong\prod_{i=1}^m A_i\), where the \(A_i\) are Artin local rings. Then for each \(i\) we have a natural surjective homomorphism (projection on the \(i\)th factor) \(\phi_i:A\to A_i\). Let \(\mathfrak a_i=\operatorname{Ker}(\phi_i)\). By \amref{1.10} the \(\mathfrak a_i\) are pairwise coprime, and \(\bigcap\mathfrak a_i=0\). Let \(\mathfrak q_i\) be the unique prime ideal of \(A_i\), and let \(\mathfrak p_i\) be its contraction \(\phi_i^{-1}(\mathfrak q_i)\). The ideal \(\mathfrak p_i\) is prime and therefore maximal by \amref{8.1}. Since \(\mathfrak q_i\) is nilpotent it follows that \(\mathfrak a_i\) is \(\mathfrak p_i\)-primary, and hence \(\bigcap\mathfrak a_i=(0)\) is a primary decomposition of the zero ideal in \(A\). Since the \(\mathfrak a_i\) are pairwise coprime, so are the \(\mathfrak p_i\), and they are therefore isolated prime ideals of \((0)\). Hence all the primary components \(\mathfrak a_i\) are isolated, and therefore uniquely determined by \(A\), by the 2nd uniqueness theorem \amref{4.11}. Hence the rings \(A_i\cong A/\mathfrak a_i\) are uniquely determined by \(A\). \proofbox

\noindent\textbf{Example.} A ring with only one prime ideal need not be Noetherian (and hence not an Artin ring). Let \(A=k[x_1,x_2,\ldots]\) be the polynomial ring in a countably infinite set of indeterminates \(x_n\) over a field \(k\), and let \(\mathfrak a\) be the ideal \((x_1,x_2^2,\ldots,x_n^n,\ldots)\). The ring \(B=A/\mathfrak a\) has only one prime ideal (namely the image of \((x_1,x_2,\ldots,x_n,\ldots)\)), hence \(B\) is a local ring of dimension 0. But \(B\) is not Noetherian, for it is not difficult to see that its prime ideal is not finitely generated.

If \(A\) is a local ring, \(\mathfrak m\) its maximal ideal, \(k=A/\mathfrak m\) its residue field, the \(A\)-module \(\mathfrak m/\mathfrak m^2\) is annihilated by \(\mathfrak m\) and therefore has the structure of a \(k\)-vector space. If \(\mathfrak m\) is finitely generated (e.g., if \(A\) is Noetherian), the images in \(\mathfrak m/\mathfrak m^2\) of a set of generators of \(\mathfrak m\) will span \(\mathfrak m/\mathfrak m^2\) as a vector space, and therefore \(\dim_k(\mathfrak m/\mathfrak m^2)\) is finite. (See \amref{2.8}.)

\noindent\textbf{Proposition 8.8.\amtarget{8.8}} \emph{Let \(A\) be an Artin local ring. Then the following are equivalent:}

\emph{i) every ideal in \(A\) is principal;}

\emph{ii) the maximal ideal \(\mathfrak m\) is principal;}

\emph{iii) \(\dim_k(\mathfrak m/\mathfrak m^2)\leq1\).}

\noindent\emph{Proof.} i) \(\Rightarrow\) ii) \(\Rightarrow\) iii) is clear.

iii) \(\Rightarrow\) i): If \(\dim_k(\mathfrak m/\mathfrak m^2)=0\), then \(\mathfrak m=\mathfrak m^2\), hence \(\mathfrak m=0\) by Nakayama's lemma \amref{2.6}, and therefore \(A\) is a field and there is nothing to prove.

If \(\dim_k(\mathfrak m/\mathfrak m^2)=1\), then \(\mathfrak m\) is a principal ideal by \amref{2.8} (take \(M=\mathfrak m\) there), say \(\mathfrak m=(x)\). Let \(\mathfrak a\) be an ideal of \(A\), other than \((0)\) or \((1)\). We have \(\mathfrak m=\mathfrak N\), hence \(\mathfrak m\) is nilpotent by \amref{8.4} and therefore there exists an integer \(r\) such that \(\mathfrak a\subseteq\mathfrak m^r\), \(\mathfrak a\not\subseteq\mathfrak m^{r+1}\); hence there exists \(y\in\mathfrak a\) such that \(y=ax^r\), \(y\notin(x^{r+1})\); consequently \(a\notin(x)\) and \(a\) is a unit in \(A\). Hence \(x^r\in\mathfrak a\), therefore \(\mathfrak m^r=(x^r)\subseteq\mathfrak a\) and hence \(\mathfrak a=\mathfrak m^r=(x^r)\). Hence \(\mathfrak a\) is principal. \proofbox

\noindent\textbf{Example.} The rings \(Z/(p^n)\) (\(p\) prime), \(k[x]/(f^n)\) (\(f\) irreducible) satisfy the conditions of \amref{8.7}. On the other hand, the Artin local ring \(k[x^2,x^3]/(x^4)\) does not: here \(\mathfrak m\) is generated by \(x^2\) and \(x^3\) (mod \(x^4\)), so that \(\mathfrak m^2=0\) and \(\dim(\mathfrak m/\mathfrak m^2)=2\).

\subsection{Exercises}

\noindent\textbf{1.} Let \(\mathfrak q_1\cap\cdots\cap\mathfrak q_n=0\) be a minimal primary decomposition of the zero ideal in a Noetherian ring, and let \(\mathfrak q_i\) be \(\mathfrak p_i\)-primary. Let \(\mathfrak p_i^{(r)}\) be the \(r\)th symbolic power of \(\mathfrak p_i\) (Chapter 4, Exercise 13). Show that for each \(i=1,\ldots,n\) there exists an integer \(r_i\) such that \(\mathfrak p_i^{(r_i)}\subseteq\mathfrak q_i\).

Suppose \(\mathfrak q_i\) is an isolated primary component. Then \(A_{\mathfrak p_i}\) is an Artin local ring, hence if \(\mathfrak m_i\) is its maximal ideal we have \(\mathfrak m_i^r=0\) for all sufficiently large \(r\), hence \(\mathfrak q_i=\mathfrak p_i^{(r)}\) for all large \(r\).

If \(\mathfrak q_i\) is an embedded primary component, then \(A_{\mathfrak p_i}\) is not Artinian, hence the powers \(\mathfrak m_i^r\) are all distinct, and so the \(\mathfrak p_i^{(r)}\) are all distinct. Hence in the given primary decomposition we can replace \(\mathfrak q_i\) by any of the infinite set of \(\mathfrak p_i\)-primary ideals \(\mathfrak p_i^{(r)}\) where \(r\geq r_i\), and so there are infinitely many minimal primary decompositions of 0 which differ only in the \(\mathfrak p_i\)-component.

\noindent\textbf{2.} Let \(A\) be a Noetherian ring. Prove that the following are equivalent:

i) \(A\) is Artinian;

ii) \(\operatorname{Spec}(A)\) is discrete and finite;

iii) \(\operatorname{Spec}(A)\) is discrete.

\noindent\amstar\textbf{3.} Let \(k\) be a field and \(A\) a finitely generated \(k\)-algebra. Prove that the following are equivalent:

i) \(A\) is Artinian;

ii) \(A\) is a finite \(k\)-algebra.

[To prove that i) \(\Rightarrow\) ii), use \amref{8.7} to reduce to the case where \(A\) is an Artin local ring. By the Nullstellensatz, the residue field of \(A\) is a finite extension of \(k\). Now use the fact that \(A\) is of finite length as an \(A\)-module. To prove ii) \(\Rightarrow\) i), observe that the ideals of \(A\) are \(k\)-vector subspaces and therefore satisfy d.c.c.]

\noindent\amstar\textbf{4.} Let \(f:A\to B\) be a ring homomorphism of finite type. Consider the following statements:

i) \(f\) is finite;

ii) the fibres of \(f^*\) are discrete subspaces of \(\operatorname{Spec}(B)\);

iii) for each prime ideal \(\mathfrak p\) of \(A\), the ring \(B\otimes_A k(\mathfrak p)\) is a finite \(k(\mathfrak p)\)-algebra (\(k(\mathfrak p)\) is the residue field of \(A_\mathfrak p\));

iv) the fibres of \(f^*\) are finite.

Prove that i) \(\Rightarrow\) ii) \(\Longleftrightarrow\) iii) \(\Rightarrow\) iv). [Use Exercises 2 and 3.]

If \(f\) is integral and the fibres of \(f^*\) are finite, is \(f\) necessarily finite?

\noindent\textbf{5.} In Chapter 5, Exercise 16, show that \(X\) is a finite covering of \(L\) (i.e., the number of points of \(X\) lying over a given point of \(L\) is finite and bounded).

\noindent\textbf{6.} Let \(A\) be a Noetherian ring and \(\mathfrak q\) a \(\mathfrak p\)-primary ideal in \(A\). Consider chains of primary ideals from \(\mathfrak q\) to \(\mathfrak p\). Show that all such chains are of finite bounded length, and that all maximal chains have the same length.

\section{Discrete Valuation Rings and Dedekind Domains}

As we have indicated before, algebraic number theory is one of the historical sources of commutative algebra. In this chapter we specialize down to the case of interest in number theory, namely to Dedekind domains. We deduce the unique factorization of ideals in Dedekind domains from the general primary decomposition theorems. Although a direct approach is of course possible one obtains more insight our way into the precise context of number theory in commutative algebra. Another important class of Dedekind domains occurs in connection with non-singular algebraic curves. In fact the geometrical picture of the Dedekind condition is: non-singular of dimension one.

The last chapter dealt with Noetherian rings of dimension 0. Here we start by considering the next simplest case, namely Noetherian \emph{integral domains} of dimension one: i.e., Noetherian domains in which every non-zero prime ideal is maximal. The first result is that in such a ring we have a unique factorization theorem for ideals:

\noindent\textbf{Proposition 9.1.\amtarget{9.1}} \emph{Let \(A\) be a Noetherian domain of dimension 1. Then every non-zero ideal \(\mathfrak a\) in \(A\) can be uniquely expressed as a product of primary ideals whose radicals are all distinct.}

\noindent\emph{Proof.} Since \(A\) is Noetherian, \(\mathfrak a\) has a minimal primary decomposition \(\mathfrak a=\bigcap_{i=1}^n\mathfrak q_i\) by \amref{7.13}, where each \(\mathfrak q_i\) is say \(\mathfrak p_i\)-primary. Since \(\dim A=1\) and \(A\) is an integral domain, each non-zero prime ideal of \(A\) is maximal, hence the \(\mathfrak p_i\) are distinct maximal ideals (since \(\mathfrak p_i\supseteq\mathfrak q_i\supseteq\mathfrak a\neq0\)), and are therefore pairwise coprime. Hence by \amref{1.16} the \(\mathfrak q_i\) are pairwise coprime and therefore by \amref{1.10} we have \(\prod\mathfrak q_i=\bigcap\mathfrak q_i\). Hence \(\mathfrak a=\prod\mathfrak q_i\).

Conversely, if \(\mathfrak a=\prod\mathfrak q_i\), the same argument shows that \(\mathfrak a=\bigcap\mathfrak q_i\); this is a minimal primary decomposition of \(\mathfrak a\), in which each \(\mathfrak q_i\) is an isolated primary component, and is therefore unique by \amref{4.11}. \proofbox

Let \(A\) be a Noetherian domain of dimension one in which every primary ideal is a prime power. By \amref{9.1}, in such a ring we shall have unique factorization of non-zero ideals into products of prime ideals. If we localize \(A\) with respect to a non-zero prime ideal \(\mathfrak p\) we get a local ring \(A_\mathfrak p\) satisfying the same conditions as \(A\), and therefore in \(A_\mathfrak p\) every non-zero ideal is a power of the maximal ideal. Such local rings can be characterized in other ways.

\subsection{Discrete valuation rings}

Let \(K\) be a field. A \emph{discrete valuation} on \(K\) is a mapping \(v\) of \(K^*\) onto \(Z\) (where \(K^*=K-\{0\}\) is the multiplicative group of \(K\)) such that

1) \(v(xy)=v(x)+v(y)\), i.e., \(v\) is a homomorphism;

2) \(v(x+y)\geq\min(v(x),v(y))\).

The set consisting of 0 and all \(x\in K^*\) such that \(v(x)\geq0\) is a ring, called the \emph{valuation ring} of \(v\). It is a valuation ring of the field \(K\). It is sometimes convenient to extend \(v\) to the whole of \(K\) by putting \(v(0)=+\infty\).

\noindent\textbf{Examples.} The two standard examples are:

1) \(K=Q\). Take a fixed prime \(p\), then any non zero \(x\in Q\) can be written uniquely in the form \(p^ay\), where \(a\in Z\) and both numerator and denominator of \(y\) are prime to \(p\). Define \(v_p(x)\) to be \(a\). The valuation ring of \(v_p\) is the local ring \(Z_{(p)}\).

2) \(K=k(x)\), where \(k\) is a field and \(x\) an indeterminate. Take a fixed irreducible polynomial \(f\in k[x]\) and define \(v_f\) just as in 1). The valuation ring of \(v_f\) is then the local ring of \(k[x]\) with respect to the prime ideal \((f)\).

An integral domain \(A\) is a \emph{discrete valuation ring} if there is a discrete valuation \(v\) of its field of fractions \(K\) such that \(A\) is the valuation ring of \(v\). By \amref{5.18}, \(A\) is a local ring, and its maximal ideal \(\mathfrak m\) is the set of all \(x\in K\) such that \(v(x)>0\).

If two elements \(x,y\) of \(A\) have the same value, that is if \(v(x)=v(y)\), then \(v(xy^{-1})=0\) and therefore \(u=xy^{-1}\) is a unit in \(A\). Hence \((x)=(y)\).

If \(\mathfrak a\neq0\) is an ideal in \(A\), there is a least integer \(k\) such that \(v(x)=k\) for some \(x\in\mathfrak a\). It follows that \(\mathfrak a\) contains every \(y\in A\) with \(v(y)\geq k\), and therefore the only ideals \(\neq0\) in \(A\) are the ideals \(\mathfrak m_k=\{y\in A:v(y)\geq k\}\). These form a single chain \(\mathfrak m\supset\mathfrak m_2\supset\mathfrak m_3\supset\cdots\), and therefore \(A\) is Noetherian.

Moreover, since \(v:K^*\to Z\) is surjective, there exists \(x\in\mathfrak m\) such that \(v(x)=1\), and then \(\mathfrak m=(x)\), and \(\mathfrak m_k=(x^k)\) \((k\geq1)\). Hence \(\mathfrak m\) is the only non-zero prime ideal of \(A\), and \(A\) is thus a Noetherian local domain of dimension one in which every non-zero ideal is a power of the maximal ideal. In fact many of these properties are characteristic of discrete valuation rings.

\noindent\textbf{Proposition 9.2.\amtarget{9.2}} \emph{Let \(A\) be a Noetherian local domain of dimension one, \(\mathfrak m\) its maximal ideal, \(k=A/\mathfrak m\) its residue field. Then the following are equivalent:}

\emph{i) \(A\) is a discrete valuation ring;}

\emph{ii) \(A\) is integrally closed;}

\emph{iii) \(\mathfrak m\) is a principal ideal;}

\emph{iv) \(\dim_k(\mathfrak m/\mathfrak m^2)=1\);}

\emph{v) Every non-zero ideal is a power of \(\mathfrak m\);}

\emph{vi) There exists \(x\in A\) such that every non-zero ideal is of the form \((x^k)\), \(k\geq0\).}

\noindent\emph{Proof.} Before we start going the rounds, we make two remarks:

(A) If \(\mathfrak a\) is an ideal \(\neq0,(1)\), then \(\mathfrak a\) is \(\mathfrak m\)-primary and \(\mathfrak a\supseteq\mathfrak m^n\) for some \(n\). For \(r(\mathfrak a)=\mathfrak m\), since \(\mathfrak m\) is the only non-zero prime ideal; now use \amref{7.16}.

(B) \(\mathfrak m^n\neq\mathfrak m^{n+1}\) for all \(n\geq0\). This follows from \amref{8.6}.

i) \(\Rightarrow\) ii) by \amref{5.18}.

ii) \(\Rightarrow\) iii). Let \(a\in\mathfrak m\) and \(a\neq0\). By remark (A) there exists an integer \(n\) such that \(\mathfrak m^n\subseteq(a)\), \(\mathfrak m^{n-1}\not\subseteq(a)\). Choose \(b\in\mathfrak m^{n-1}\) and \(b\notin(a)\), and let \(x=a/b\in K\), the field of fractions of \(A\). We have \(x^{-1}\notin A\) (since \(b\notin(a)\)), hence \(x^{-1}\) is not integral over \(A\), and therefore by \amref{5.1} we have \(x^{-1}\mathfrak m\not\subseteq\mathfrak m\) (for if \(x^{-1}\mathfrak m\subseteq\mathfrak m\), \(\mathfrak m\) would be a faithful \(A[x^{-1}]\)-module, finitely generated as an \(A\)-module). But \(x^{-1}\mathfrak m\subseteq A\) by construction of \(x\), hence \(x^{-1}\mathfrak m=A\) and therefore \(\mathfrak m=Ax=(x)\).

iii) \(\Rightarrow\) iv). By \amref{2.8} we have \(\dim_k(\mathfrak m/\mathfrak m^2)\leq1\), and by remark (B) \(\mathfrak m/\mathfrak m^2\neq0\).

iv) \(\Rightarrow\) v). Let \(\mathfrak a\) be an ideal \(\neq(0),(1)\). By remark (A) we have \(\mathfrak a\supseteq\mathfrak m^n\) for some \(n\); from \amref{8.8} (applied to \(A/\mathfrak m^n\)) it follows that \(\mathfrak a\) is a power of \(\mathfrak m\).

v) \(\Rightarrow\) vi). By remark (B), \(\mathfrak m\neq\mathfrak m^2\), hence there exists \(x\in\mathfrak m\), \(x\notin\mathfrak m^2\). But \((x)=\mathfrak m^r\) by hypothesis, hence \(r=1\), \((x)=\mathfrak m\), \((x^k)=\mathfrak m^k\).

vi) \(\Rightarrow\) i). Clearly \((x)=\mathfrak m\), hence \((x^k)\neq(x^{k+1})\) by remark (B). Hence if \(a\) is any non-zero element of \(A\), we have \((a)=(x^k)\) for exactly one value of \(k\). Define \(v(a)=k\) and extend \(v\) to \(K^*\) by defining \(v(ab^{-1})=v(a)-v(b)\). Check that \(v\) is well-defined and is a discrete valuation, and that \(A\) is the valuation ring of \(v\). \proofbox

\subsection{Dedekind domains}

\noindent\textbf{Theorem 9.3.\amtarget{9.3}} \emph{Let \(A\) be a Noetherian domain of dimension one. Then the following are equivalent:}

\emph{i) \(A\) is integrally closed;}

\emph{ii) Every primary ideal in \(A\) is a prime power;}

\emph{iii) Every local ring \(A_\mathfrak p\) (\(\mathfrak p\neq0\)) is a discrete valuation ring.}

\noindent\emph{Proof.} i) \(\Longleftrightarrow\) iii) by \amref{9.2} and \amref{5.13}.

ii) \(\Longleftrightarrow\) iii). Use \amref{9.2} and the fact that primary ideals and powers of ideals behave well under localization: \amref{4.8}, \amref{3.11}. \proofbox

A ring satisfying the conditions of \amref{9.3} is called a \emph{Dedekind domain}.

\noindent\textbf{Corollary 9.4.\amtarget{9.4}} \emph{In a Dedekind domain every non-zero ideal has a unique factorization as a product of prime ideals.}

\noindent\emph{Proof.} \amref{9.1} and \amref{9.3}. \proofbox
